%%%%%%%%%%%%%%%%%%%%%%%%%%%%%%%%%%%%%%%%%%%%%%%%%%%%%%%%%%%%%%%%%%%%%%%%%%%%%%%%%%%%%%
\section{Correctness Arguments}
\label{sec:guarantees}
%%%%%%%%%%%%%%%%%%%%%%%%%%%%%%%%%%%%%%%%%%%%%%%%%%%%%%%%%%%%%%%%%%%%%%%%%%%%%%%%%%%%%%
%=====================================================================================
\subsection{\primitive implementation}
\label{sec:log-correctness}
%=====================================================================================
In this section we argue that \primitive implementation in \Cref{alg:bbca-pt1} satisfies \primitive properties.
%-------------------------------------------------------------------------------------
\subsubsection{Integrity}
\label{sec:bbca-integrity}
%-------------------------------------------------------------------------------------
Integrity is satisfied 
by correct validators only sending a \echo message for an \initiate message signed by the dedicate sender of the instance
(see \Cref{alg:bbca-pt1}, line \ref{ln:integrity}).

%-------------------------------------------------------------------------------------
\subsubsection{No duplication}
\label{sec:bbca-no-duplication}
%-------------------------------------------------------------------------------------
No duplication trivial to satisfy by a correct validator checking that it has not \primitivesmall-delivered for the same instance before (see \Cref{alg:bbca-pt1}, lines \ref{ln:no-duplication-1}, \ref{ln:no-duplication-2}).
%-------------------------------------------------------------------------------------
\subsubsection{Consistency}
\label{sec:bbca-consistency}
%-------------------------------------------------------------------------------------
Let's assume by contradiction that there exist correct validators $p$ and $q$ who \primitivesmall-deliver $(id, sn, b)$ and $id, sn', b'$ with $sn\neq sn'$ or $b\neq b'$.

There are three conditions which, upon being satisfied, cause a validator to trigger $\bbcadeliver(p, id, sn, b)$: (a) by receiving $2f+1$ matching \ready messages before a timeout for instance $id$ or a $\bbcaabort(sn)$ event, (b) because of a $\brbdeliver$ event, and (c) by receiving $f+1$ matching \checkpoint messages. 

We exhaustively examine all the possible combinations.
\begin{enumerate}
	\item Both $p$ and $q$ trigger \bbcadeliver\ by condition (a).
	\item Both $p$ and $q$ trigger \bbcadeliver\ by condition (b).
	\item One of $p$ and $q$ triggers \bbcadeliver\ by condition (a) and one of $p$ and $q$ triggers \bbcadeliver\ by condition (b).
	\item Both $p$ and $q$ trigger \bbcadeliver\ by condition (c).
	\item One of $p$ and $q$ triggers \bbcadeliver\ by conditions (a) or (b) and one of $p$ and $q$ triggers \bbcadeliver\ by condition (c).
	\item One of $p$ and $q$ triggers \bbcadeliver\ before a timeout on instance $id$ or \bbcaabort\ for $sn$ or $sn'$ has been triggered and the other after.
\end{enumerate}

\paragraph{Case 1.} In this case, $p$ has received $2f+1$ \echo messages for $(id, sn, b)$ and $q$ has received $2f+1$ \echo messages for $(id, sn', b')$, which suggests that there exists a correct validator who sent and \echo message for both triplets.
A contradiction, since a correct validator sends only one \echo message per instance (see \Cref{alg:bbca-pt1}, line \ref{ln:integrity}).

\paragraph{Case 2.}
A contradiction due to the consistency property of \brb (no two different correct validators trigger \brbdeliver for different messages) .

\paragraph{Case 3.}
Let's assume, without loss of generality, that $p$ \primitivesmall-delivers before a timeout on instance $id$ or a \bbcaabort\ event (condition (a))  and $q$ after (condition (b)).
Therefore $p$ has received $2f+1$ \echo messages for $(id, sn, b)$.
This suggests that the external validity predicate for \brb is satisfied: there exist some ready certificate $C_R\neq\emptyset$ for $id, sn', b'$, i.e., $2f+1$ validators sent an \echo message for $id, sn', b'$.
A contradiction, since a correct validator sends only one \echo message per instance.

\paragraph{Case 4.}
In this case, $p$ has received $f+1$ \checkpoint messages from $f+1$ validators for $(id, sn, b)$, among which there exists at least one correct validator, let's call it $p1$.
Similarly, $q$ has received $f+1$ \checkpoint messages from $f+1$ validators for $(id, sn', b')$, among which there exists at least one correct validator, let's call it $q1$.
If $p1=q1$, then $p1$ has triggered \bbcadeliver for different triplets, a contradiction.
If $p1\neq p2$, the have triggered \bbcadeliver due to either (a) or (b). Hence, we fall back to cases $1,2$, and $3$.

\paragraph{Case 5.}
Let's assume, without loss of generality, that $p$ has received $f+1$ \checkpoint messages from $f+1$ validators for $(id, sn, b)$,  among which there exists at least one correct validator, let's call it $p1$.
In this case $p1$ has triggered \bbcadeliver due to either (a) or (b).
Therefore, we fall back to cases $1,2$, and $3$ for the correct validators $p1$ and $q$.

%-------------------------------------------------------------------------------------
\subsubsection{Unique sequencing}
\label{sec:bbca-sequencing}
%-------------------------------------------------------------------------------------
Again, let's assume by contradiction that the there exist correct validators $p$ and $q$ who \primitivesmall deliver $(id, sn, b)$ and $id', sn, b'$ with $id\neq id'$ and $b\neq b'$.
Since $sn\neq\emptyslot$, both $p$ and $q$  trigger \bbcadeliver\ before a timeout on instance $id$ or $id$ and before \bbcaabort\ for $sn$ .

Therefore, $p$ has received $2f+1$ \echo messages for $(id, sn, b)$ and $q$ has  has received $2f+1$ \echo messages for $(id', sn, b')$, which suggests that there exists a correct validator who sent and \echo message for both triplets.
A contradiction, since a correct validator sends only one \echo message per slot number (see \Cref{alg:bbca-pt1}, line \ref{ln:integrity}).

%-------------------------------------------------------------------------------------
\subsubsection{Validity}
\label{sec:bbca-validity}
%-------------------------------------------------------------------------------------
\paragraph{Synchronous network communication.} If the network is synchronous, i.e. we are after GST, no correct validator triggers a \bbcaabort\ event and there is no invocation of \bbcacast\ with the same slot number $s$, by the weak sequencing validity property, every correct validator delivers the same triplet $(id,sn,b)$.

Let's assume now that we are before GST.

\paragraph{Timed-out instances.} 
There exist two mutually exclusive cases: (a) at least $f+1$ correct validators who timed-out for instance $id$ and (b) at most $f$ validators timed-out for instance $id$.

In case (a), all correct validators will receive a \abort message and, therefore, send an \abort message themselves.
Therefore, the sender of the instance will receive $2f+1$ \abort messages and invoke $\brbcast(id, sn', b)$ where  $sn'=sn$ or $sn'=\emptyslot$, depending on the ready certificates in the \abort messages.
Since the sender is correct, by \brb validity, all correct validators will \brb-deliver $(id,sn',b)$, and, therefore, all correct validators who had timed-out before sending an \abort message \primitivesmall-deliver $id, sn', b$.
Finally, by \primitive consistency, all correct validators will \primitivesmall-deliver $id, sn', b$.

In case (b), there exist at least $f+1$ correct validators who \primitivesmall-deliver $(id,sn,b)$.
They all deliver $(id,sn,b)$ by \primitive consistency property. 
Therefore, all correct validators, who had not delivered before, will eventually receive $f+1$ \checkpoint messages and \primitivesmall-deliver $(id,sn,b)$.

\paragraph{Conflicting instances.} Let us now assume that some other sender invokes $\bbcacast(id',sn,b')$.
If some correct validator sends an \echo message for $(id',sn',b')$ before receiving an \initiate message for $(id,sn,b)$, then 
it is possible that some correct validators will not gather $2f+1$ \echo messages for $(id,sn,b)$, resulting in a timeout for instance $id$. 
Then validity follows by the same arguments as in timed-out instances paragraph.

\paragraph{Finalizations.} If some correct validator triggers a \bbcaabort\ event before it \primitivesmall-delivers, it will timeout for instance $id$.
 Then validity follows by the same arguments as in timed-out instances paragraph.

%-------------------------------------------------------------------------------------
\subsubsection{Weak sequencing validity}
\label{sec:bbca-weak-seq-validity}
%-------------------------------------------------------------------------------------
Since there is no other instance \primitivesmall-bcast  with $sn$, no correct validator has sent an $\langle \echo, id', sn, b'\rangle$ message.
Also, since the sender is correct, it invokes \bbcacast for $id,sn,b$ such that $\mathcal{P}(id,sn,b)=\true$. 
Moreover, since we are after GST no correct validator triggers $\bbcaabort(sn)$ and no timeout occurs for $id$.
Therefore, every correct validator sends an $\langle \echo, id, sn, b\rangle$ message and, thus, every correct validator receives $2f+1$  valid \echo messages and sends a \ready message.
Similarly, every correct validator receives $2f+1$ valid \ready messages and \primitivesmall-delivers $(id,sn,b)$ 
Finally, by consistency and unique sequencing properties, all validators \primitivesmall-deliver $(id,sn,b)$.

%-------------------------------------------------------------------------------------
\subsubsection{External validity}
\label{sec:consistency}
%-------------------------------------------------------------------------------------
If $P(b)$ is not satisfied, no correct validator sends an \echo message.
Therefore, no correct validator can deliver $(id,sn,b)$.

%-------------------------------------------------------------------------------------
\subsubsection{Finalization}
\label{sec:bbca-finalization}
%-------------------------------------------------------------------------------------
A $\bbcaabort(sn)$ event trivially triggers a \bbcaannounce\ event by the implementation (see \Cref{alg:bbca-pt1}, line \ref{ln:announcement}).
Let's now argue about the second part of the property.
In our implementation the proof $\pi$ is a ready certificate.
Predicate $\mathcal{Q}$ corresponds to no ready certificate being empty.
Therefore, we have $2f+1$ empty ready certificates, among which at least $f+1$ are from correct validators.
Since $f+1$ correct validators did not sent a \ready message, there exist at most $2f$ validators which sent a \ready message for slot number $sn$.
Therefore, no correct validator could have \primitivesmall-delivered $id, sn, b$ for any block $b$ and any instance $id$.


%=====================================================================================
\subsection{\system implementation correctness}
\label{sec:log-correctness}
%=====================================================================================

Let us first formally specify the properties we need \system to guarantee.

\begin{itemize}
	\item \textbf{Safety}
	\begin{itemize}
	\item \textit{Consistency}: If two correct validators $p$ and $q$ commit blocks $b$ and $b'$ respectively to the same slot, then $b$=$b'$.
	\end{itemize}
	\item \textbf{Liveness}
	\begin{itemize}
	\item \textit{Infinite log growth}: Eventually all correct validators commit a value (a block or the special $\bot$ value) for any slot.
	\item \textit{Eventual progress}: After GST the log grows with non $\bot$ values.
	\end{itemize}
\end{itemize}

We argue that the implementation in \Cref{alg:sb-impl-pt1} satisfies the properties of \system.

%-------------------------------------------------------------------------------------
\subsubsection{Consistency}
\label{sec:consistency}
%-------------------------------------------------------------------------------------
We prove consistency by contradiction%
\footnote{Slotless blocks inherit consistency from the blocks with slot numbers by which they are committed.}.
Let us assume that validators $p$ and $q$ commit blocks $p$ and $q$ at slot $sn$ with $b\neq b'$.
There exist six potential ways in which the blocks have been committed.
\begin{enumerate}
	\item Both $p$ and $q$ \primitivesmall-delivered $b$ and $b'$ with $sn$.
	\item One of $p$, $q$ \primitivesmall-delivered the corresponding block and the other decided via the fallback protocol.
	\item Both $p$, $q$ decided via the fallback protocol.
	\item One of $p$ and $q$ \primitivesmall-delivered $b$ and the other pulled $b'$ from another validator.
	\item One of $p$ and $q$ decided $b$ via the fallback protocol and the other pulled $b'$ form another validator.
	\item Both $p$, $q$ pulled $b$ and $b'$ from other validators.
\end{enumerate}

In the first case, $p$ and $q$ \primitivesmall-delivered $(id, sn, b$) and $(id', sn, b')$ respectively.
Then, from \primitive unique sequencing property ($\primitiveproperty4$), $b=b'$, a contradiction.

In the second case, let us assume without loss of generality that $p$ \primitivesmall-delivered $id, sn, b$ and $q$ decided $sn,b'$.
Therefore, within \primitive, $p$ received $2f+1$ \ready messages, $f+1$ out of which from a set of correct validators $S_1$.
Let $S_2$ be the set of validators who invoked the \bbcaabort\ events leading to $q$ deciding.
$S_1, S_2$ intersect in at least one correct validator; let's name the correct validator $p_c$.
By belonging in set $S_1$, $p_c$ sent a \ready message for block $b$,
and as argued above already, $b$ is unique with such a ready message for $sn$.
Hence, $q$ must decide on $b$, because $S_2$ includes the ready-certificate $C_{R}$ by $p_c$.
A contradiction.

In the third case, let us assume that $p$ and $q$ decide  $sn, b$ and $sn, b'$ respectively.
By the \vbc safety, no two correct processes decide differently.
Hence, $b=b'$. A contradiction. 

In the fourth case, let's assume without loss of generality that $q$ pulled $b'$ form another validator.
This means that $b'$ was referenced in $f+1$ blocks the $q$ \primitivesmall-delivered such each of them was \primitivesmall-cast by a different validator.
Among the $f+1$ validators at least one is correct. Let's name the correct validator $p_c$.
Therefore, $p_c$ \primitivesmall-delivered $b'$ with $sn$. 
A contradiction follows from the arguments for case (1) for $p$ and $p_c$.
Similarly, case (5), assuming that $p$ decided via the fallback consensus, reduces to case (2) for $p$ and $p_c$.
Finally, in case (6), both blocks $b$ and $b'$ are referenced by $f+1$ blocks.
Therefore, there exist correct validators $p_c$ and $q_c$ who \primitivesmall-delivered $b$ and $b'$ and the contradiction follows from case (1).


%-------------------------------------------------------------------------------------
\subsubsection{Infinite log growth}
\label{sec:growth}
%-------------------------------------------------------------------------------------
\lemma{\textit{Totality}: If a correct validator \primitivesmall-delivers a block $(id,sn,b)$ or decides $sn$ for $b$ via the fallback consensus protocol, then every correct validator either \primitivesmall-delivers $(id,sn,b)$ or decides $b$ for slot number $sn$ via the fallback consensus protocol.}
\label{lemma:totality}
\proof{
	If the \primitive sender $p_s$ for instance $id$ is correct, by \primitive integrity, $p_s$ has invoked \primitivesmall-cast for $(id,sn,b)$ and, by \primitive validity, eventually every correct validator delivers $(id,sn,b)$.
	We examine the case where \primitive sender is not correct.
	There are two possible cases; either $sn=\emptyslot$ or $sn\neq\emptyslot$.
	In the first case, $b$ is delivered in the second phase of \primitive, i.e. via a validated reliable broadcast instance which guarantees totality.
	In the second case, if $f+1$ correct validators \primitivesmall-deliver $(id,sn,b)$, then every correct validator \primitivesmall-delivers $(id,sn,b)$ via the \primitive checkpointing mechanism.
	Otherwise, for the \primitive-ledger protocol, $sn$ will cause a timeout for at least $f+1$ correct validators. 
	This is because, by \primitive consistency property, no correct validator will deliver any block other than $b$ for $sn$.
	Therefore, at least $f+1$ correct validators will \primitivesmall-cast a block with the \finalize flag which, by \primitive validity, will be \primitivesmall-delivered by all correct validators.
	Therefore, all correct validators to \primitivesmall-cast a block with the \finalize flag and eventually trigger \vbc for $sn$.
	By \vbc liveness and following the argumentation of \Cref{sec:consistency}, all correct validators will eventually decide $b$ for $sn$.	 
	Finally, if a correct validator decides $sn$ for $b$ via the fallback consensus protocol, then by the properties of the fallback consesus, all correct validators eventually decide $sn$ for $b$.
}

\lemma{\textit{Causality}: If a correct validator $p$ either \primitivesmall-delivers $(id,sn,b)$ or decides $b$ for slot number $sn$ via the fallback consensus protocol, where $b\neq\bot$, then $p$ eventually finalizes $b$ with $sn$.}
\label{lemma:causality}
\proof{
	 Let us assume by contradiction that there exists a causal predecessor $b'$ of $b$ which $p$ never delivers, and therefore $p$ cannot finalize $b$.
	 In the first case, where $p$ either \primitivesmall-delivers, by \primitive external validity property, the external predicate $P$ is valid for $b$, i.e., $p$ has \primitivesmall-delivered $b'$. A contradiction.
	 In the second case, where $p$ decides $b$ via the fallback finalization protocol, since $b\neq\bot$, there exists at least one non empty ready certificate for $b$.
	 This suggests that $2f+1$ validators sent \echo messages for $b$, and therefore, they had already delivered $b'$.
	 Among them, there exist at least $f+1$ correct. 
	 Therefore, by \Cref{lemma:totality}, all correct validators, including $p$, will eventually deliver $p$.
	 A contradiction.
}

Using \Cref{lemma:totality} and \Cref{lemma:causality} we prove the infinite log property for $\snmin$, the first slot number which has not been committed.
Then by a simple induction argument we can show that the property holds for any slot number.

Let's assume by contradiction that some correct validator $p$ never commits a value in slot $\snmin$.
This will eventually trigger a timeout for slot number $\snmin$ and $p$ will trigger $\bbcaabort(\snmin)$.
Let $t$ be the number of correct validators who \primitivesmall-deliver a block $b$ for slot number $\snmin$.
We examine two mutually exclusive sub-cases.
\begin{enumerate}
	\item $t\ge\ f+1$: At least $f+1$ correct validators \primitivesmall-deliver a block $b$ for slot number $\snmin$ and $b\neq\bot$.
	\item $t<f$: There exist at least $f$ more correct validators whose timer for $\snmin$ expires.
\end{enumerate}

In the first case, since the $f+1$ correct validators are correct, by \Cref{lemma:totality}, all correct validators, including $p$, will either \primitivesmall-deliver $b$ with $sn$ or \vbc-decide $b$ for $sn$ and by \Cref{lemma:causality}, $p$ will finalize slot $sn$ with $b$.
Since, by assumption, all slots before $\snmin$ are committed, then $\snmin$ also gets committed.
 A contradiction.


In the second case $f+1$ correct validators will \primitivesmall-bcast a block with a \finalize flag for $\snmin$.
By \primitive validity, all correct validators will \primitivesmall-deliver the blocks, causing eventually all correct validators to \primitivesmall-bcast a block with a \finalize flag for $\snmin$.
Therefore, all correct validators will \primitivesmall-deliver at least $2f+1$ blocks with a \finalize flag for $\snmin$ carrying a valid ready certificate  and will invoke
$\bcpropose([\snmin, b,\mathcal{C}],\mathcal{Q})$,
where $\mathcal{C}$ be set of the ready certificates of the first $2f+1$ blocks which satisfy the predicate $\mathcal{Q}$ for $\snmin$ and some block $b$, possibly also $\bot$.
By \vbc liveness, all correct validators, including $p$, receive $\vbtobdecide([\snmin, b,\mathcal{C}])$.
By \Cref{lemma:causality}, $p$ finalizes $b$ with $\snmin$.
Since $p$ has committed all slots before $\snmin$, it also commits slot $\snmin$.
A contradiction.

%-------------------------------------------------------------------------------------
\subsubsection{Eventual progress}
\label{sec:growth}
%-------------------------------------------------------------------------------------
If a correct validator $p$ invokes \primitive for some block $b$ which includes a set of causal references to a set of blocks $B$, such that $p$ has either \primitivesmall-delivered or \vbc-decided for each block in $B$, then, by \Cref{lemma:totality}, all correct validators will \primitivesmall-deliver
	of \vbc-decide for all blocks in $B$.
Therefore, the external validity predicate for \primitive will be satisfied for all correct validators.
Moreover, by the weak sequencing validity property of \primitive ($\primitiveproperty6$), after GST, if some correct validator $p$ invokes \bbcacast\ for some slot number $sn$, all correct validators will \bbcadeliver\ $sn$ with some block $b\neq\bot$, and hence add it to the log, if the following conditions hold: 
\begin{enumerate}
	\item No correct validator invokes \bbcaabort for $sn$.
	\item There exists no other validator who invokes \primitive for $sn$.
\end{enumerate}

The first condition is satisfied, since we are after GST and, therefore no validator will timeout.
The second condition depends on the ticket master implementation.
As we cover in \Cref{sec:ticketmaster}, the ticketmaster functionality ensures that all correct validators are assigned slot numbers and, moreover, slot numbers are assigned exclusively.
The round robin strategy satisfies this by assigning slot numbers to validators in a deterministic way.
In the strategies where the ticket master is a server, we achieve this by rotating the ticket master among validators, so that eventually some correct validator assumes the role.

